
\section{Introduction}

Retrieval-augmented generation (RAG) has become a prominent approach for building user-facing NLP applications, such as systems for question answering (QA), fact-checking, and customer support~\cite{petroni-etal-2021-kilt, wang2019superglue}. Typically, a RAG system consists of a retriever and a downstream language model (LM). Given a user question, the retriever finds relevant passages from a corpus 
and the LM uses these passages to generate a response. This formulation admits a multitude of choices: 
what retrieval model to use, %
how to divide the documents into retrieval chunks, %
and how to prompt %
or finetune %
the LM to use the retrieved information, to name only a few of the simplest design decisions.

The best design for a RAG system is not necessarily universal across data domains, corpus sizes, and cost/latency budgets. To tune their own RAG systems, practitioners traditionally need hand annotations for test questions, passages to retrieve (to assess the retriever), and responses to generate, labeled specifically for their target domain. Alternatively, they may evaluate different approaches in production by collecting human preferences that compare the candidate systems. Unfortunately, both of these strategies demand high expertise and impose considerable annotation costs.

Model-based evaluation is
an inexpensive
strategy to test generative output quality \cite{zheng2023judging}. For instance, the open-source RAGAS framework ~\cite{ragas2023} prompts an LM for evaluating the \textit{relevance} of retrieved information and the \textit{faithfulness} and \textit{accuracy} of generated responses. Unfortunately, such strategies currently rely for evaluation on a fixed set of heuristically hand-written prompts, offering little adaptability to various evaluation contexts and no guarantees about quality. %






To evaluate RAG systems rapidly and accurately, we propose ARES, the \textbf{A}utomated \textbf{R}AG \textbf{E}valuation \textbf{S}ystem.
ARES is the first automated RAG evaluation system to generate tailored LLM judges for each component of a RAG pipeline, leading to substantial boosts in evaluation precision and accuracy compared to existing approaches like RAGAS.
Furthermore, unlike existing RAG evaluation systems, ARES provides 
confidence intervals for its scoring
by leveraging prediction-powered inference (PPI; \citealt{angelopoulos2023predictionpowered}).
Given a corpus of documents and a RAG system, ARES reports three evaluation scores: context relevance (is the retrieved information pertinent to the test question), answer faithfulness (is the response generated by the language model properly grounded in the retrieved context), and answer relevance (is the response also relevant to the question). A good RAG system finds relevant contexts and generates answers that are both faithful and relevant.

Many existing RAG evaluation frameworks require substantial human annotations for scoring.
ARES significantly improves data efficiency during evaluation by only requiring three inputs: an in-domain passage set, a human preference validation set of approximately 150 annotated datapoints or more, and few-shot examples of in-domain queries and answers (e.g. five examples or more), which are used for prompting LLMs in synthetic data generation.

Given the corpus of in-domain passages, ARES proceeds in three stages. First, it leverages an LM to construct a synthetic dataset of question--answer pairs, derived from the passages in the corpus. 
Second, it defines three separate judge models to perform three classification tasks (context relevance, answer faithfulness, and answer relevance). These judges are lightweight models fine-tuned against a contrastive learning objective. 
Third, ARES scores the different RAG systems being assessed using prediction-powered inference (PPI; \citealt{angelopoulos2023predictionpowered}) to improve model-based evaluation accuracy and provide statistical confidence intervals for RAG scoring.
PPI utilizes a small set of human annotated datapoints for computing its confidence intervals; we designate this annotated set as our \textit{human preference validation set}, which is composed of approximately 150 annotated datapoints or more that designate both positive and negative examples for context relevance, answer faithfulness, and answer relevance. 





We conduct extensive empirical evaluations, demonstrating that ARES accurately scores RAG systems across the six knowledge-intensive datasets in KILT and SuperGLUE, beating existing automated evaluation approaches like RAGAS by 59.3 and 14.4 percentage points on average across context relevance and answer relevance evaluation accuracy, respectively. 
Additionally, ARES accurately calculates answer hallucination occurrences in the AIS attribution dataset \cite{rashkin2022measuring}, predicting within 2.5 percentage points of the ground truth average for answer hallucinations.
Compared to annotation-based evaluation methods, ARES is substantially more accurate and efficient, requiring 78\% less annotations than the baseline approach.
We also find that ARES consistently distinguishes competitive RAG systems that are only a few points apart in ground-truth metrics. This precision enables ARES to guide the development and comparison of competitive approaches and configurations.

We make the ARES code and datasets publicly available \href{https://github.com/stanford-futuredata/ARES}{on Github}.















